% 実コードを読んでみよう（第7章と第8章の橋渡し）
\section{実コードを読んでみよう}

実践に入る前に、リポジトリに実装されている最もシンプルなモジュール「LedModule」の実コードを読んでみましょう。ここまで学んできたパターンが、実際のコードでどう使われているかを確認します。

%% ── Q1: LedModule.h ──────────────────────────
\begin{qabox}{LedModule.h はどうなってるの？}

まずはヘッダーファイル（\texttt{lib/LedModule/LedModule.h}）です。

\begin{lstlisting}
// LedModule.h — LED制御モジュール
#pragma once
#include <Arduino.h>
#include "IModule.h"

// --- Config構造体 ---
struct LedConfig {
    uint8_t ledPin;  // LED出力ピン
};

// --- Data構造体 ---
struct LedData {
    bool state = false;
    uint8_t brightness = 0;
};

// SystemDataの前方宣言はIModule.h内で行われている

// --- モジュール実装 ---
class LedModule : public IModule {
private:
    LedConfig _config;

public:
    LedModule(const LedConfig& config);
    bool init() override;
    void updateOutput(SystemData& data) override;
};
\end{lstlisting}

ここまで学んだパターンがすべて揃っています。

\begin{itemize}
  \item \textbf{Config構造体}（第3章） — \texttt{LedConfig} にピン番号をまとめている
  \item \textbf{Data構造体}（第3章） — \texttt{LedData} にデフォルト値が設定されている
  \item \textbf{IModule継承}（第4章） — \texttt{: public IModule} で統一インターフェースを実装
  \item \textbf{前方宣言}（第5章） — \texttt{SystemData.h} をインクルードしていない（IModule.hの中で前方宣言済み）
  \item \textbf{override}（第4章） — 親クラスの仮想関数を上書きすることを明示
  \item \textbf{updateOutput}（第5章） — LEDは出力モジュールなので \texttt{updateOutput} をオーバーライド
\end{itemize}

\begin{notebox}[ヘッダーには宣言だけ]
メソッドの中身（実装）は書かれていません。「こういう関数がありますよ」という宣言だけです。実装は \texttt{.cpp} ファイルに書きます。これが第7章で説明した「\texttt{.h} + \texttt{.cpp}」の分離ルールです。
\end{notebox}

\end{qabox}

%% ── Q2: LedModule.cpp ──────────────────────────
\begin{qabox}{LedModule.cpp はどうなってるの？}

次に実装ファイル（\texttt{lib/LedModule/LedModule.cpp}）です。

\begin{lstlisting}
// LedModule.cpp — LED制御モジュールの実装
#include "LedModule.h"
#include "SystemData.h"
#include <Arduino.h>

LedModule::LedModule(const LedConfig& config) : _config(config) {}

bool LedModule::init() {
    pinMode(_config.ledPin, OUTPUT);
    Serial.println("[Led] init OK");
    return true;
}

void LedModule::updateOutput(SystemData& data) {
    digitalWrite(_config.ledPin, data.led.state ? HIGH : LOW);
}
\end{lstlisting}

たった16行です。各メソッドを見ていきましょう。

\subsection{コンストラクタ}

\begin{lstlisting}
LedModule::LedModule(const LedConfig& config) : _config(config) {}
\end{lstlisting}

Configを受け取って \texttt{\_config} にコピーするだけです。ピン番号はハードコードせず、Configから取得します。

\subsection{init()}

\begin{lstlisting}
bool LedModule::init() {
    pinMode(_config.ledPin, OUTPUT);
    Serial.println("[Led] init OK");
    return true;
}
\end{lstlisting}

\texttt{\_config.ledPin} でピン番号にアクセスしています。ログフォーマットは \texttt{[モジュール名] メッセージ} の統一形式です。成功時は \texttt{true} を返します。

\subsection{updateOutput()}

\begin{lstlisting}
void LedModule::updateOutput(SystemData& data) {
    digitalWrite(_config.ledPin, data.led.state ? HIGH : LOW);
}
\end{lstlisting}

SystemDataから \texttt{data.led.state} を読み取り、LEDのON/OFFを制御しています。LED自身は「なぜ光るのか」を知りません。ロジックフェーズが \texttt{data.led.state} を書き換え、LEDモジュールはそれに従うだけです。これが3フェーズモデルの「出力モジュール」の役割です。メソッド名が \texttt{updateOutput} なので、このモジュールが出力専用であることが一目で分かります。

\begin{notebox}[\texttt{.cpp} では SystemData.h をインクルード]
\texttt{.h} では前方宣言だけでしたが、\texttt{.cpp} では \texttt{\#include "SystemData.h"} でフル定義を取得しています。\texttt{data.led.state} のようにメンバにアクセスするには、SystemDataの中身を知る必要があるためです。
\end{notebox}

\end{qabox}

%% ── Q3: 全体の流れ ──────────────────────────
\begin{qabox}{このモジュールはどうやって動くの？}

LedModuleが動くまでの全体の流れを確認しましょう。

\begin{enumerate}
  \item \texttt{ProjectConfig.h} で設定を定義:
\begin{lstlisting}
const LedConfig LED_CONFIG = { .ledPin = 2 };
\end{lstlisting}

  \item \texttt{main.cpp} でインスタンスを作成し、配列に登録:
\begin{lstlisting}
LedModule ledModule(LED_CONFIG);
IModule* outputModules[] = { &ledModule };
\end{lstlisting}

  \item \texttt{setup()} で \texttt{init()} が呼ばれ、ピンが初期化される
  \item \texttt{loop()} の出力フェーズで \texttt{updateOutput(systemData)} が呼ばれる
  \item ロジックフェーズが \texttt{systemData.led.state} を書き換えると、次のループでLEDに反映される
\end{enumerate}

これが「Config → インスタンス化 → init → updateInput/updateOutput」の基本サイクルです。次の章では、このサイクルをゼロから自分で作ります。

\end{qabox}
